\documentclass[10pt]{article}
\usepackage{style/blueprintdoc}
\geometry{landscape,margin=11mm,top=14mm,bottom=12mm}
\setdoccode{NNV-CS}\setdoctitle{Cheat sheet — neural network verification}
\setlength{\columnsep}{7mm}
\begin{document}
{\bpdisplay\bfseries\LARGE\color{bpInk}Neural network verification — cheat sheet}\quad{\bpmono\small\color{bpGraphite}NNV-CS · REV A · Sept 2026}\par\vspace{1mm}{\color{bpRule}\hrule height 0.5pt}\vspace{2mm}
\begin{multicols}{3}\footnotesize
\subsection*{Vocabulary}
\textbf{Property} — ``for every input $x$ in region $R$: $\varphi(x)$''. \textbf{Region} — the inputs the claim covers; a design decision. \textbf{Proof} — holds for \emph{every} point. \textbf{Counterexample} — one concrete failing input, in your units. \textbf{Timeout} — an honest ``too hard''. \textbf{Sound} — never claims a false property. \textbf{Complete} — always answers, given time. \textbf{Robustness ($\varepsilon$-ball)} — same class for every input within $\varepsilon$. \textbf{Embedding gap} — requirement in feet/mmHg vs solver in normalised tensors. \textbf{Hyperproperty} — about two evaluations (monotonicity, fairness): needs an encoding. \textbf{ONNX / VNN-LIB} — network / property exchange formats. \textbf{VNN-COMP} — annual competition (2025: 8 tools).

\subsection*{Vehicle syntax card}
\begin{vclbox}[types \& declarations]
type Input  = Tensor Real [7]
type Output = Tensor Real [3]
low = 0          -- name an index
@network
f : Input -> Output
@parameter
eps : Real
@parameter(infer=True)
n : Nat
@dataset
xs : Vector Input n
@dataset
ys : Vector (Index 3) n
\end{vclbox}
\begin{vclbox}[predicates \& properties]
valid : Input -> Bool
valid x = 0 <= x ! 0 <= 100 and x ! 5 == 0

advises : Input -> Index 3 -> Bool
advises x c = forall j . j != c =>
                 f x ! c > f x ! j

@property
p : Bool
p = forall x . valid x and x ! 1 <= 91
      => not (advises x low)

robustAround : Input -> Index 3 -> Bool
robustAround x c = forall d .
  (forall i . -eps <= d ! i <= eps)
  and valid (x + d) => advises (x + d) c

@property
q : Vector Bool n
q = foreach i . robustAround (xs ! i) (ys ! i)
\end{vclbox}
Annotations (\texttt{@network}, \texttt{@parameter}, \texttt{@dataset}) go on their own line; put named indices \emph{above} the network. Records (\texttt{x.spo2}) read well but do not combine with datasets in 0.27.1.

\subsection*{Command line}
\begin{shellbox}[verify / typecheck / export]
(*@\tprompt@*)vehicle typecheck -s spec.vcl
(*@\tprompt@*)vehicle verify -s spec.vcl \
   -n f:model.onnx -y p \
   -d xs:data.idx -d ys:labels.idx \
   -p eps:0.5 --solver Marabou \
   -a "--timeout=30" [-c cache/]
(*@\tprompt@*)vehicle compile queries \
   --format VNNLibQueries -o out/ \
   -s spec.vcl -n f:model.onnx -e p
\end{shellbox}
\texttt{-y} may repeat; omit it to verify all properties. \texttt{--solver} takes a name on PATH or a file path. Vehicle has no timeout of its own — always pass Marabou's, and cap the wall clock too.

\subsection*{Reading results}
\texttt{result: 🗸 - Marabou proved no counterexample exists} → \verified\\[0.5mm]
\texttt{result: ✗ - Marabou found a counterexample} then \texttt{x: [ … ]} → \counterexample\ (units of the spec)\\[0.5mm]
\texttt{result: ? - Marabou timed out} → \unknown\\[0.5mm]
Vector properties end with \texttt{verified: a/n  falsified: b/n  timed-out: c/n}.\\[0.5mm]
Warning about strict inequalities: \texttt{>} is relaxed to \texttt{>=}; ties on the boundary count as failures (and can appear as spurious counterexamples in fairness specs — add a margin).

\subsection*{Limits to remember}
\begin{itemize}
  \item Quantified variables only \textbf{linearly} (no $x\cdot y$).
  \item No \textbf{alternating} quantifiers (\texttt{forall x . exists y}).
  \item \textbf{One} network application per query: hyperproperties need a twin (self-composed) network — and two \emph{identical} copies can make the solver loop; prove those structurally.
  \item No Boolean inputs: pin flags with \texttt{==}.
  \item \texttt{if}/\texttt{or}/\texttt{max} multiply queries; prefer relations.
  \item Marabou's ONNX reader: Gemm/MatMul/Add/Relu/Reshape/… — no Concat, Slice, Div. Fold normalisation into layer 1.
  \item NP-complete: 16–32 hidden units → seconds; vision-scale → minutes to never.
\end{itemize}

\subsection*{What a good property looks like}
\begin{enumerate}
  \item One sentence about \emph{every} input in a named region, in the problem's units.
  \item The region is plausible and bounded on every input.
  \item Thresholds sit away from where the training labels change (here: a full unit).
  \item You know what a counterexample would look like and who must see it.
  \item You have written down what the property does \emph{not} cover.
\end{enumerate}
\end{multicols}
\end{document}
